\section*{अध्याय \ref{ch:b1:poly} --- बहुपद}

\begin{solution}{exo:b1:poly:1}
$X^5 - 1 = (X^2 + X + 1)(X^3 - X^2 + 1) + (-X - 2)$। पद: पहले $X^3 B$ घटाइए,
फिर $-X^2 B$, फिर $B$; शेषफल $-X - 2$ की घात $1 < 2$ है। \emph{$X = 1$ पर
जाँच:} $\;0 = 3 \times 1 + (-3)$।

$2X^4 + X^3 - X + 3 = (X^2 - 2)(2X^2 + X + 4) + (X + 11)$। \emph{$X = 0$ पर
जाँच:} $\;3 = (-2)(4) + 11$।
\end{solution}

\begin{solution}{exo:b1:poly:2}
$j = \eu^{2\iu\pi/3}$, $j^3 = 1$ के साथ $X^2 + X + 1 = (X - j)(X - j^2)$। वह
$Q_n = X^{2n} + X^n + 1$ को तभी \omterm{def:b1:arith:divides}{विभाजित करता है} जब $j$ और $j^2$ $Q_n$ के मूल
हों; और चूँकि $Q_n$ के गुणांक वास्तविक हैं, $Q_n(j^2) = \conj{Q_n(j)}$, अतः
शर्त केवल $Q_n(j) = 0$ है। अब $Q_n(j) = j^{2n} + j^n + 1$ $3$ के सापेक्ष $n$
पर निर्भर करता है:
\begin{itemize}
  \item $n \equiv 0$: $Q_n(j) = 1 + 1 + 1 = 3 \neq 0$;
  \item $n \equiv 1$: $Q_n(j) = j^2 + j + 1 = 0$;
  \item $n \equiv 2$: $Q_n(j) = j^4 + j^2 + 1 = j + j^2 + 1 = 0$.
\end{itemize}
अतः $X^2 + X + 1 \mid X^{2n} + X^n + 1$ ठीक तब जब $3 \nmid n$।
\end{solution}

\begin{solution}{exo:b1:poly:3}
\cref{prop:b1:poly:multiplicity} से $(X-1)^2 \mid P$ तभी जब $P(1) = P'(1) =
0$:
\[
P(1) = 2 + a + b = 0, \qquad P'(1) = 4 + 3a + 2b = 0 .
\]
हल करने पर: $b = -a - 2$ और $4 + 3a - 2a - 4 = a = 0$, अतः $a = 0$, $b =
-2$: $P = X^4 - 2X^2 + 1 = (X^2 - 1)^2 = (X-1)^2 (X+1)^2$, जो वास्तविक
गुणनखंडन है।
\end{solution}

\begin{solution}{exo:b1:poly:4}
$\C$ पर $X^3 - 1 = (X - 1)(X - j)(X - j^2)$ ($j = \eu^{2\iu\pi/3}$), और $\R$
पर $(X - 1)(X^2 + X + 1)$।

$\R$ पर $X^4 + X^2 + 1 = (X^2 + X + 1)(X^2 - X + 1)$ (गुणा करके देखिए, अथवा
$X^4 + X^2 + 1 = (X^2+1)^2 - X^2$ पर ध्यान दीजिए); $\C$ पर प्रत्येक द्विघात
टूट जाता है: मूल $j, j^2$ और $-j, -j^2$, अर्थात् $\eu^{\pm 2\iu\pi/3},
\eu^{\pm\iu\pi/3}$।

$\C$ पर $X^6 - 1 = \prod_{k=0}^{5} (X - \eu^{\iu k\pi/3})$, और $\R$ पर:
\[
X^6 - 1 = (X-1)(X+1)(X^2 + X + 1)(X^2 - X + 1),
\]
जहाँ संयुग्मी जोड़े $\eu^{\pm 2\iu\pi/3}$ और $\eu^{\pm \iu\pi/3}$ समूहबद्ध
किए गए हैं।
\end{solution}

\begin{solution}{exo:b1:poly:5}
\begin{enumerate}
  \item मान लीजिए $p/q$ (न्यूनतम पदों में) \omterm{def:b1:poly:def}{इकाई-अग्र} पूर्णांक \omterm{def:b1:poly:def}{बहुपद} $X^3 +
  \dots + c_0$ का मूल है: $P(p/q) = 0$ में हर हटाने पर $p^3 =
  -q\,(\text{पूर्णांक})$ मिलता है, अतः $q \mid p^3$; सहअभाज्यता $q = \pm 1$
  पर बाध्य कर देती है: अर्थात् मूल पूर्णांक $p$ है, और $p \mid c_0$ ($c_0$
  अलग कीजिए)। यहाँ उम्मीदवार $6$ को विभाजित करते हैं: जाँचने पर $P(1) = 0$,
  $P(2) = 0$, $P(3) = 0$। अतः $P = (X-1)(X-2)(X-3)$।
  \item व्येता: $s_1 = 6$, $s_2 = 11$, $s_3 = 6$। वर्गों का योग: $s_1^2 -
  2s_2 = 36 - 22 = 14 = 1 + 4 + 9$, जैसा अपेक्षित था। प्रतिलोमों का योग:
  $\frac{s_2}{s_3} = \frac{11}{6} = 1 + \frac12 + \frac13$, जैसा अपेक्षित
  था।
\end{enumerate}
\end{solution}

\begin{solution}{exo:b1:poly:6}
\omterm{met:b1:arith:euclid}{यूक्लिडीय कलनविधि} का पहला भाग-पद:
\[
(X + 1)(X^3 - X^2 + X - 1)
= X^4 - X^3 + X^2 - X + X^3 - X^2 + X - 1 = X^4 - 1 ,
\]
अतः $X^4 - 1$ का $X^3 - X^2 + X - 1$ से भाग यथार्थ है (भागफल $X + 1$, शेषफल
$0$), और कलनविधि तुरंत रुक जाती है:
\[
\gcd(X^4 - 1,\; X^3 - X^2 + X - 1) = X^3 - X^2 + X - 1
\]
(जो पहले से \omterm{def:b1:poly:def}{इकाई-अग्र} है)। बेज़ू संबंध तुच्छ है: $\gcd = 0 \cdot (X^4 - 1) +
1 \cdot (X^3 - X^2 + X - 1)$। गुणनखंडन से संगति की जाँच: $X^3 - X^2 + X - 1
= (X - 1)(X^2 + 1)$, जो सचमुच $X^4 - 1 = (X-1)(X+1)(X^2+1)$ के उभयनिष्ठ
अखंडनीय गुणनखंडों का गुणनफल है।
\end{solution}

\begin{solution}{exo:b1:poly:7}
चूँकि $\R$ पर $P \geq 0$, उसके वास्तविक मूलों की \omterm{def:b1:poly:derivative}{बहुलता} सम है (विषम \omterm{def:b1:poly:derivative}{बहुलता}
वाले मूल पर $P$ चिह्न बदल देता)। \cref{cor:b1:poly:factorization} का प्रयोग
करके और युग्मन करके लिखिए
\[
P = c \prod_i (X - a_i)^{2k_i} \prod_j \bigl((X - z_j)(X - \conj
z_j)\bigr)^{n_j},
\]
जहाँ $c > 0$ ($+\infty$ पर व्यवहार से)। मान लीजिए
\[
S = \sqrt c\, \prod_i (X - a_i)^{k_i} \prod_j (X - z_j)^{n_j}
\in \C[X],
\]
जिससे $P = S\,\conj S$, जहाँ $\conj S$ के गुणांक संयुग्मी हैं। $A, B \in
\R[X]$ के साथ $S = A + \iu B$ तोड़िए: तब
\[
P = (A + \iu B)(A - \iu B) = A^2 + B^2 .
\]
\end{solution}

\begin{solution}{exo:b1:poly:8}
बिंदुओं $0, 1, 2$ के साथ लाग्रांज (\cref{thm:b1:poly:lagrange}):
\[
P = 1\cdot\frac{(X-1)(X-2)}{(0-1)(0-2)} + 3\cdot\frac{X(X-2)}{1\cdot(1-2)}
+ 2\cdot\frac{X(X-1)}{2\cdot 1}
= \frac{(X-1)(X-2)}{2} - 3X(X-2) + X(X-1).
\]
प्रसार करने पर: $\frac{X^2 - 3X + 2}{2} - 3X^2 + 6X + X^2 - X =
-\frac{3}{2}X^2 + \frac{7}{2}X + 1$।

निकाय: $c = 1$ के साथ $P = aX^2 + bX + c$; $a + b + 1 = 3$; $4a + 2b + 1 =
2$। तीसरे में से दूसरे का दुगुना घटाने पर $2a - 1 = -4$, अतः $a = -\frac32$,
$b = \frac72$। वही \omterm{def:b1:poly:def}{बहुपद}: $P = -\frac32 X^2 + \frac72 X + 1$। (जाँच $P(2) =
-6 + 7 + 1 = 2$।)
\end{solution}

\begin{solution}{exo:b1:poly:9}
$P = X^{2n+1} - 1$: $P' = (2n+1)X^{2n} \geq 0$, अतः \omterm{def:b1:poly:def}{बहुपद} फलन वर्धमान है
($0$ को छोड़कर कठोरतः), जिसकी सीमाएँ $\mp\infty$ हैं: वह $\R$ पर ठीक एक बार
लुप्त होता है ($x = 1$ पर)।

मान लीजिए $E_n = \sum_{k=0}^{n} \frac{X^k}{k!}$। तब $E_n' = E_{n-1} = E_n -
\frac{X^n}{n!}$। कोई बहु मूल $a$ $E_n(a) = E_n'(a) = 0$
(\cref{prop:b1:poly:multiplicity}) संतुष्ट करता, अतः $\frac{a^n}{n!} =
E_n(a) - E_n'(a) = 0$, इसलिए $a = 0$; पर $E_n(0) = 1 \neq 0$। कोई बहु मूल
नहीं।
\end{solution}

\begin{solution}{exo:b1:poly:10}
\begin{enumerate}
  \item आगमन (दोनों आधार स्थितियाँ सत्य हैं)। $\cos(n+1)\theta +
  \cos(n-1)\theta = 2\cos\theta\cos n\theta$ का प्रयोग करते हुए:
        \[
        T_{n+1}(\cos\theta) = 2\cos\theta \cos n\theta -
        \cos(n-1)\theta = \cos(n+1)\theta .
        \]
  \item $T_n(\cos\theta) = 0$ तभी जब $n\theta \equiv \frac\pi2 \pmod \pi$:
  संख्याएँ
        \[
        x_k = \cos\Bigl(\frac{(2k+1)\pi}{2n}\Bigr),
        \qquad k = 0, 1, \dots, n-1,
        \]
        $\intoo{-1}{1}$ के $n$ भिन्न बिंदु हैं (कोण $\intoo{0}{\pi}$ में
        हैं, जहाँ $\cos$ \omterm{def:b1:logic:inj}{एकैकी} है), और वे सब $T_n$ के मूल हैं; चूँकि $\deg
        T_n = n$ (पुनरावृत्ति से, जिसका अग्र गुणांक आगमन से $n \geq 1$ के
        लिए $2^{n-1}$ है), अतः ये \emph{सभी} मूल हैं, और हर एक सरल है।
  \item $x = \cos\theta \in \intcc{-1}{1}$ के लिए: $\abs{T_n(x)} = \abs{\cos
  n\theta} \leq 1$, जहाँ समता तभी जब $n\theta \equiv 0 \pmod\pi$, अर्थात्
  $n+1$ बिंदुओं $y_k = \cos\frac{k\pi}{n}$, $k = 0, \dots, n$ पर, जहाँ
  $T_n(y_k) = (-1)^k$। (यही समदोलन $2^{1-n}T_n$ को $\intcc{-1}{1}$ पर लघुतम
  उच्चतम-मानक वाला घात-$n$ का \omterm{def:b1:poly:def}{इकाई-अग्र बहुपद} बनाता है --- उपपत्ति इस अध्याय
  की सप्ताहांत समस्या में।)
\end{enumerate}
\end{solution}

\begin{solution}{exo:b1:poly:11}
$P = c\prod_i (X - a_i)^{m_i}$ लिखिए। गुणनफल नियम (कई गुणनखंडों तक
विस्तारित) देता है
\[
P' = c\sum_{i} m_i (X - a_i)^{m_i - 1} \prod_{k \neq i} (X -
a_k)^{m_k},
\]
और $P$ से भाग देने पर: $\frac{P'}{P} = \sum_i \frac{m_i}{X - a_i}$ (परिमेय
फलनों के रूप में, अर्थात् मूलों से दूर)।

मान लीजिए $w$ $P'$ का मूल है। यदि $w$ $a_i$ में से एक है, तो वह तुच्छ रूप से
उत्तल आवरण में है। अन्यथा $w$ पर मान लेने पर:
\[
0 = \sum_i \frac{m_i}{w - a_i}
= \sum_i m_i\, \frac{\conj w - \conj a_i}{\abs{w - a_i}^2} .
\]
संयुग्मी लेने पर: $\sum_i \lambda_i (w - a_i) = 0$, जहाँ $\lambda_i =
\frac{m_i}{\abs{w - a_i}^2} > 0$। अतः
\[
w = \frac{\sum_i \lambda_i a_i}{\sum_i \lambda_i} :
\]
जो मूलों $a_i$ का उत्तल संयोजन है (धनात्मक भार, जिनका योग मानकीकरण के बाद
$1$ है)। अतः $P'$ का प्रत्येक मूल $P$ के मूलों के उत्तल आवरण में है।
\end{solution}

\begin{solution}{exo:b1:poly:12}
इकाई के तीनों घनमूलों पर $(1 + X)^n$ के मानों का योग कीजिए:
\[
2^n + (1 + j)^n + (1 + j^2)^n
= \sum_{k=0}^n \binom nk\,\bigl(1 + j^k + j^{2k}\bigr)
= 3\sum_{k\,:\,3\mid k}\binom nk ,
\]
क्योंकि $1 + j^k + j^{2k}$ ऐसा गुणोत्तर योग है जो $3 \mid k$ होने पर $3$ के
बराबर है और अन्यथा $\frac{j^{3k} - 1}{j^k - 1} = 0$ के। अब $1 + j = \frac12
+ \iu\frac{\sqrt3}2 = \eu^{\iu\pi/3}$ और $1 + j^2 = \conj{1 + j} =
\eu^{-\iu\pi/3}$, अतः $(1+j)^n + (1+j^2)^n = 2\cos\frac{n\pi}3$ और
\[
\sum_{k\geq0}\binom n{3k} = \frac{2^n + 2\cos\frac{n\pi}3}{3} .
\]
जाँच: $n = 3$: $\frac{8 + 2\cos\pi}3 = 2 = \binom30 + \binom33$; $n = 6$:
$\frac{64 + 2}3 = 22 = 1 + 20 + 1$।
\end{solution}

\begin{solution}{pb:b1:poly:1}
\textbf{1.} $T_2 = 2X^2 - 1$; $T_3 = 2X(2X^2 - 1) - X = 4X^3 - 3X$; $T_4 =
2X\,T_3 - T_2 = 8X^4 - 8X^2 + 1$; $T_5 = 2X\,T_4 - T_3 = 16X^5 - 20X^3 +
5X$। सर्वसमिका $T_3(\cos\theta) = \cos3\theta$ ठीक \cref{ex:b1:complex:cos3}
का $\cos3\theta = 4\cos^3\theta - 3\cos\theta$ ही है।

\textbf{2.} $n = 1, 2$ के लिए सत्य। यदि $T_{n-1}$, $T_n$ की घातें $n-1$, $n$
और अग्र गुणांक $2^{n-2}$, $2^{n-1}$ हैं, तो $2X\,T_n$ की घात $n+1$ और अग्र
गुणांक $2^n$ है, जबकि $T_{n-1}$ की घात कम है: अतः $T_{n+1}$ की घात $n + 1$
और अग्र गुणांक $2^n$ है। सम-विषमता: यदि $T_{n-1}$ की सम-विषमता $n - 1$ जैसी
है और $T_n$ की $n$ जैसी, तो $2X\,T_n$ और $T_{n-1}$ दोनों की सम-विषमता $n +
1$ जैसी है, अतः $T_{n+1}$ की भी।

\textbf{3.} यदि सभी $\theta$ के लिए $P(\cos\theta) = \cos n\theta$, तो $P$
और $T_n$ $\intcc{-1}1$ के प्रत्येक बिंदु पर मेल खाते हैं --- और वह अनंत
\omterm{def:b1:logic:sets}{समुच्चय} है --- अतः $P - T_n$ के अनंत मूल हैं और वह शून्य \omterm{def:b1:poly:def}{बहुपद} है
(\cref{cor:b1:poly:nroots})।

\textbf{4.} $x = \cos\theta$ के लिए: $T_m(T_n(\cos\theta)) = T_m(\cos
n\theta) = \cos(mn\theta) = T_{mn}(\cos\theta)$, और $2T_mT_n(\cos\theta) =
2\cos m\theta\cos n\theta = \cos(m+n)\theta + \cos\abs{m - n}\theta$। दोनों
सर्वसमिकाएँ $\intcc{-1}1$ पर सत्य हैं, अतः प्रश्न 3 के तर्क से बहुपदीय
सर्वसमिकाओं के रूप में भी।

\textbf{5.} $x_k$ $n$ भिन्न सरल मूल हैं और अग्र गुणांक $2^{n-1}$ है:
\[
T_n = 2^{n-1}\prod_{k=0}^{n-1}
\Bigl(X - \cos\frac{(2k+1)\pi}{2n}\Bigr) .
\]
एक-दूसरे के बीच आना: कोण $0 < \frac{\pi}{2n} < \frac\pi n < \frac{3\pi}{2n}
< \frac{2\pi}n < \dots < \pi$ $y$-कोणों $\frac{k\pi}n$ और $x$-कोणों
$\frac{(2k+1)\pi}{2n}$ के बीच एकांतरित होते हैं; और चूँकि $\cos$
$\intcc0\pi$ पर कठोरतः ह्रासमान है, मान उलटे क्रम में एक-दूसरे के बीच आ जाते
हैं: $y_n < x_{n-1} < y_{n-1} < \dots < x_0 < y_0$। दो क्रमागत चरम बिंदुओं
के बीच ठीक एक मूल बैठता है, जैसा $\cos n\theta$ का चित्र संकेत देता है।

\textbf{6.} $2\cosh a\cosh b = \cosh(a + b) + \cosh(a - b)$
(\cref{prop:b1:functions:hyprules}) के साथ आगमन: $T_{n+1}(\cosh t) = 2\cosh
t\cosh nt - \cosh(n-1)t = \cosh(n+1)t$। $x \geq 1$ के लिए $t \geq 0$ के साथ
$x = \cosh t$ लिखिए; तब $\eu^t = x + \sqrt{x^2 - 1}$ और $\eu^{-t} = x -
\sqrt{x^2 - 1}$, अतः
\[
T_n(x) = \cosh(nt)
= \frac{(x + \sqrt{x^2-1})^n + (x - \sqrt{x^2-1})^n}2 .
\]
$x > 1$ के लिए पहला पद $\frac12(1)^n$ को कठोरतः पार कर जाता है और गुणोत्तर
रूप से बढ़ता है: $T_n(x) > 1$।

\textbf{7.} दे मॉयवर: $\cos n\theta = \Re\bigl((\cos\theta +
\iu\sin\theta)^n\bigr) = \sum_{2j \leq n}\binom n{2j}
\cos^{n-2j}\theta\,(\iu\sin\theta)^{2j}$, और $(\iu\sin\theta)^{2j} =
(-\sin^2\theta)^j = (\cos^2\theta - 1)^j$। $x = \cos\theta$ प्रतिस्थापित
करके और प्रश्न 3 का आह्वान करके:
\[
T_n(x) = \sum_{0\leq 2j\leq n}\binom n{2j}x^{n-2j}(x^2 - 1)^j .
\]
$n = 3$ के लिए: $\binom30 x^3 + \binom32 x(x^2 - 1) = x^3 + 3x^3 - 3x = 4x^3
- 3x$, जैसा प्रश्न 1 में है।

\textbf{8.} $T_n(1) = \cos(n\cdot0) = 1$; $T_n(-1) = \cos(n\pi) = (-1)^n$;
$T_n(0) = \cos\frac{n\pi}2$, जो विषम $n$ के लिए $0$ है और सम $n$ के लिए
$(-1)^{n/2}$।

\textbf{9.} $U_n(\cos\theta) = \frac{\sin(n+1)\theta}{\sin\theta}$ के लिए
आगमन: $U_0 = 1$ और $U_1 = 2X$ के लिए सत्य ($\sin2\theta =
2\sin\theta\cos\theta$); और पद योग-से-गुणनफल सर्वसमिका $\sin(n+2)\theta =
2\cos\theta\, \sin(n+1)\theta - \sin n\theta$ है। अब $\theta$ में
$T_n(\cos\theta) = \cos n\theta$ का अवकलन कीजिए:
$-\sin\theta\,T_n'(\cos\theta) = -n\sin n\theta$, अतः $\theta \notin \pi\Z$
के लिए:
\[
T_n'(\cos\theta) = n\,\frac{\sin n\theta}{\sin\theta}
= n\,U_{n-1}(\cos\theta) ,
\]
और \omterm{def:b1:poly:def}{बहुपद} $T_n'$ तथा $nU_{n-1}$, जो $\intoo{-1}1$ पर मेल खाते हैं, बराबर हैं।

\textbf{10.} $\abs{\sin(n+1)\theta} = \abs{\sin n\theta\cos\theta + \cos
n\theta\sin\theta} \leq \abs{\sin n\theta} + \abs{\sin\theta}$, और आगमन
$\abs{\sin n\theta} \leq n\abs{\sin\theta}$ देता है। अतः $\intoo{-1}1$ पर
$\abs{U_{n-1}} \leq n$ और वहीं $\abs{T_n'} = n\abs{U_{n-1}} \leq n^2$;
$\pm1$ पर परिबंध सीमाएँ लेकर बढ़ जाता है (अथवा सीधे: पुनरावृत्ति से
$U_{n-1}(1) = n$, आगमन से $U_n(1) = n + 1$, और सम-विषमता $U_{n-1}(-1) =
(-1)^{n-1}n$ देती है)। इस प्रकार $T_n'(1) = n^2$ और $T_n'(-1) =
(-1)^{n-1}n^2$: परिबंध $n^2$ सिरों पर प्राप्त होता है।

\textbf{11.} $\theta$ के सापेक्ष $\sin\theta\,T_n'(\cos\theta) = n\sin
n\theta$ (प्रश्न 9) का अवकलन कीजिए:
\[
\cos\theta\,T_n'(\cos\theta) - \sin^2\theta\,T_n''(\cos\theta)
= n^2\cos n\theta = n^2\,T_n(\cos\theta) .
\]
$x = \cos\theta$ और $\sin^2\theta = 1 - x^2$ के साथ: $\intcc{-1}1$ पर
$x\,T_n' - (1 - x^2)T_n'' = n^2T_n$, अतः सर्वत्र: $y = T_n$ के लिए $(1 -
x^2)y'' - xy' + n^2y = 0$। $T_2 = 2x^2 - 1$ के लिए जाँच: $(1 - x^2)(4) -
x(4x) + 4(2x^2 - 1) = 4 - 4x^2 - 4x^2 + 8x^2 - 4 = 0$।

\textbf{12.} $\widetilde T_n$ \omterm{def:b1:poly:def}{इकाई-अग्र} है (प्रश्न 2) और $\intcc{-1}1$ पर
$\abs{\widetilde T_n} = 2^{1-n}\abs{T_n} \leq 2^{1-n}$, जहाँ $n + 1$ बिंदुओं
$y_k$ (\cref{exo:b1:poly:10}) पर $\widetilde T_n(y_k) = (-1)^k2^{1-n}$: अतः
मानक ठीक $2^{1-n}$ है, जो एकांतरित चिह्नों के साथ प्राप्त होता है।

\textbf{13.} $\widetilde T_n$ और $P$ दोनों घात $n$ के \omterm{def:b1:poly:def}{इकाई-अग्र} हैं, अतः
अग्र पद कट जाते हैं: $\deg D \leq n - 1$। $y_k$ पर: $D(y_k) = (-1)^k2^{1-n}
- P(y_k)$, और $\abs{P(y_k)} \leq \norm P_\infty < 2^{1-n}$ $D(y_k)$ के चिह्न
को कठोरतः $(-1)^k2^{1-n}$ वाला होने पर बाध्य कर देता है।

\textbf{14.} प्रत्येक $k = 0, \dots, n-1$ के लिए $D$ $y_{k+1}$ और $y_k$ के
बीच चिह्न बदलता है: मध्यवर्ती मान गुणधर्म से $D$ का इन $n$ जोड़ों में
असंयुक्त खुले अंतरालों में से प्रत्येक में एक मूल है --- अर्थात् घात $\leq n
- 1$ के अशून्य \omterm{def:b1:poly:def}{बहुपद} के $n$ भिन्न मूल, जो असंभव है। और $D = 0$ भी असंभव है
(मानक भिन्न हैं)। विरोधाभास: घात $n$ के किसी \omterm{def:b1:poly:def}{इकाई-अग्र} $P$ के लिए $\norm
P_\infty < 2^{1-n}$ नहीं हो सकता, और यही चेबिशेव की प्रमेय है।

\textbf{15.} अब केवल $\abs{P(y_k)} \leq 2^{1-n}$, अतः $(-1)^k D(y_k) =
2^{1-n} - (-1)^kP(y_k) \geq 2^{1-n} - \abs{P(y_k)} \geq 0$। मान लीजिए किसी
भीतरी $y_k$ ($0 < k < n$) पर $D(y_k) = 0$: तब $P(y_k) = (-1)^k2^{1-n}$, अतः
$\abs P$ अपना उच्चतम $2^{1-n}$ भीतरी बिंदु $y_k$ पर प्राप्त करता है, जिससे
$P'(y_k) = 0$ (भीतरी चरम); और $T_n'(y_k) = nU_{n-1}(y_k) = 0$, क्योंकि
$\sin(n\cdot\frac{k\pi}n) = 0$ --- अतः $\widetilde T_n'(y_k) = 0$ भी, और
$D'(y_k) = 0$: अर्थात् $y_k$ $D$ का कम से कम \omterm{def:b1:poly:derivative}{बहुलता} $2$ वाला मूल है।

\textbf{16.} \omterm{def:b1:poly:derivative}{बहुलता} सहित $D$ के मूल गिनिए। मान लीजिए $z$ ऐसे भीतरी बिंदुओं
$y_k$ की संख्या है जिनके लिए $D(y_k) = 0$ (प्रश्न 15 से हर एक द्विक मूल), और
$e \in \{0, 1, 2\}$ ऐसे सिरों ($y_0$ या $y_n$) की संख्या जिनके लिए $D = 0$
(हर एक कम से कम सरल मूल)। जिस अंतराल $(y_{k+1}, y_k)$ के दोनों सिरों पर $D
\neq 0$ हो, वह कठोरतः एकांतरित चिह्न रखता है, अतः उसमें एक भीतरी मूल है। हर
लुप्त होने वाला भीतरी बिंदु अधिक से अधिक अपने दो पड़ोसी अंतराल बिगाड़ता है
और हर लुप्त होने वाला सिरा अधिक से अधिक एक: अतः कम से कम $n - 2z - e$ अंतराल
फिर भी एक-एक मूल देते हैं, जो $y$-मूलों से भिन्न हैं। कुल: घात $\leq n - 1$
के \omterm{def:b1:poly:def}{बहुपद} के लिए \omterm{def:b1:poly:derivative}{बहुलता} सहित कम से कम $(n - 2z - e) + 2z + e = n$ मूल: अतः $D
= 0$ और $P = \widetilde T_n$। न्यूनकारी अद्वितीय है।

\textbf{17.} ऐफ़ीन \omterm{def:b1:logic:map}{प्रतिचित्रण} $t \mapsto x = \frac{a+b}2 + \frac{b-a}2\,t$
एक एकैकी आच्छादन $\intcc{-1}1 \to \intcc ab$ है। यदि $P$ घात $n$ का
\omterm{def:b1:poly:def}{इकाई-अग्र} है, तो $Q(t) = P(x(t))$ $t$ में ऐसा \omterm{def:b1:poly:def}{बहुपद} है जिसका अग्र गुणांक
$\bigl(\frac{b-a}2\bigr)^n$ है, और $\sup_{\intcc ab}\abs P =
\sup_{\intcc{-1}1}\abs Q$। \omterm{def:b1:poly:def}{इकाई-अग्र बहुपद} $Q/\bigl(\frac{b-a}2\bigr)^n$ का
उच्चतम-मानक $\geq 2^{1-n}$ है (प्रश्न 13--14), अतः
\[
\sup_{\intcc ab}\abs P \geq \Bigl(\frac{b-a}2\Bigr)^n 2^{1-n}
= 2\Bigl(\frac{b-a}4\Bigr)^n ,
\]
जहाँ समता ठीक $P(x) = \bigl(\frac{b-a}2\bigr)^n \widetilde
T_n\bigl(t(x)\bigr)$ के लिए है (प्रश्न 16)।

\textbf{18.} $\widetilde T_3 = \frac{T_3}4 = X^3 - \frac34X$; $\widetilde
T_3{}' = 3X^2 - \frac34$ $\pm\frac12$ पर लुप्त होता है। मान: $\widetilde
T_3(-1) = -\frac14$, $\widetilde T_3(-\tfrac12) = \frac14$, $\widetilde
T_3(\tfrac12) = -\frac14$, $\widetilde T_3(1) = \frac14$: अर्थात् निरपेक्ष
मान $\frac14$ वाले चार एकांतरित चरम --- अतः $\norm{\widetilde T_3}_\infty =
\frac14$, और चेबिशेव की प्रमेय से $\intcc{-1}1$ पर किसी \omterm{def:b1:poly:def}{इकाई-अग्र} त्रिघात का
उच्चतम-मानक इससे छोटा नहीं है।

\textbf{19.} $\omega$ घात $n + 1$ का \omterm{def:b1:poly:def}{इकाई-अग्र} है, अतः चेबिशेव की प्रमेय
(घात $n+1$) से $\norm\omega_\infty \geq 2^{-n}$, जहाँ समता यदि और केवल यदि
$\omega = \widetilde T_{n+1} = 2^{-n}T_{n+1}$ (प्रश्न 16), अर्थात् यदि और
केवल यदि बिंदु $T_{n+1}$ के $n + 1$ मूल हों। चेबिशेव बिंदुओं के साथ
त्रुटि-गुणक $\norm\omega_\infty$ $2^{-n}$ होता है --- अर्थात् न्यूनतम संभव।

\textbf{20.} $5 \times 36^\circ = 180^\circ$, अतः $T_5(c) = \cos180^\circ =
-1$: $16c^5 - 20c^3 + 5c + 1 = 0$। $x = -1$ जाँचने पर: $-16 + 20 - 5 + 1 =
0$, और प्रसार करने पर पुष्टि होती है
\[
16x^5 - 20x^3 + 5x + 1 = (x + 1)\bigl(4x^2 - 2x - 1\bigr)^2 .
\]
चूँकि $c = \cos36^\circ \neq -1$, $c$ $4x^2 - 2x - 1$ का मूल है, जिसके मूल
$\frac{1 \pm \sqrt5}4$ हैं; और चूँकि $c > 0$,
\[
\cos36^\circ = \frac{1 + \sqrt5}4 .
\]
संगति: $\cos72^\circ = T_2(c) = 2c^2 - 1 = 2\cdot\frac{3 + \sqrt5}8 - 1 =
\frac{\sqrt5 - 1}4$, अर्थात् वही मान जो \cref{exo:b1:complex:8} में मिला था।

\textbf{21.} $\sqrt{1.1^2 - 1} = \sqrt{0.21} \approx 0.458$, अतः $x +
\sqrt{x^2-1} \approx 1.558$ और $(1.558)^{10} \approx 84.5$, जबकि $(1.1 -
0.458)^{10} \approx 0.01$: $T_{10}(1.1) \approx \frac{84.5 + 0.01}2 \approx
42$। जो \omterm{def:b1:poly:def}{बहुपद} उस अंतराल पर $\intcc{-1}1$ में बँधा है, वह किनारे से दसवाँ भाग
बाहर जाते ही $40$ पार कर चुका होता है: किसी खंड पर परिबद्धता उससे एक इंच
बाहर के विषय में कुछ नहीं कहती।

\textbf{22.} प्रश्न 7 के $T_p$ वाले सूत्र में पद $j = 0$ $X^p$ है; शेष हर पद
$0 < 2j < p$ वाला $\binom p{2j}$ लिए हुए है (ध्यान दीजिए $2j \neq p$,
क्योंकि $p$ विषम है), और वह \cref{thm:b1:arith:fermat} की उपपत्ति के पहले पद
से $p$ से विभाज्य है। अतः $T_p - X^p$ का प्रत्येक गुणांक $p$ का गुणज है।
जाँच: $T_3 - X^3 = 3X^3 - 3X = 3(X^3 - X)$; $T_5 - X^5 = 15X^5 - 20X^3 + 5X
= 5(3X^5 - 4X^3 + X)$।

\textbf{23.} $\intcc ab = \intcc01$ और $n = 2$ के साथ प्रश्न 17 से: लघुतम
विचलन $2\bigl(\frac14\bigr)^2 = \frac18$, जो
$\bigl(\frac12\bigr)^2\widetilde T_2(2x - 1) = \frac14\bigl((2x-1)^2 -
\frac12\bigr) = x^2 - x + \frac18$ से प्राप्त होता है। $\intcc01$ पर शून्य
के सबसे निकट रहने वाला \omterm{def:b1:poly:def}{इकाई-अग्र} द्विघात $x^2 - x + \frac18$ है, जिसका
उच्चतम-मानक $\frac18$ है।

\textbf{24.} (क) दृढ़ता --- जिस \omterm{def:b1:poly:def}{बहुपद} के मूल उसकी घात से अधिक हों वह शून्य
है --- ने अद्वितीयता के सिद्धांत (प्रश्न 3), त्रिकोणमितीय सर्वसमिकाओं के
बहुपदीय सर्वसमिकाओं में स्थानांतरण (प्रश्न 4, 7, 9, 11), और चरम प्रमेय की
उपपत्ति के दोनों मूल-गणना तर्कों (प्रश्न 14, 16) को शक्ति दी। (ख)
\cref{ch:b1:complex} की त्रिकोणमिति (दे मॉयवर, योग-से-गुणनफल) और
\cref{ch:b1:functions} के अतिपरवलयिक फलनों ने इस कुल के पीछे की हर सर्वसमिका
दी; प्रतिस्थापन $x = \cos\theta$ ही सेतु है। (ग) \cref{ch:b1:arith} से आई
\omterm{def:b1:arith:divides}{विभाज्यता} $p \mid \binom p{2j}$ ने गुणांक-सूत्र को प्रश्न 22 की \omterm{def:b1:arith:congruence}{सर्वांगसमता}
में बदल दिया।

\textbf{25.} चेबिशेव की प्रमेय एक अनंतविमीय कुल (सभी \omterm{def:b1:poly:def}{इकाई-अग्र बहुपद}) पर
इष्टतमीकरण को परिमित संयोजनशास्त्र में बदल देती है: $\widetilde T_n$ से
बेहतर कोई प्रतिद्वंद्वी उससे ऐसे न्यून-घात \omterm{def:b1:poly:def}{बहुपद} से भिन्न होता जिसे $n$ बार
चिह्न बदलना पड़ता --- अर्थात् उसकी घात से एक मूल अधिक। अतः समदोलन का
प्रतिरूप कोई कुतूहल नहीं, बल्कि इष्टतमता का प्रमाणपत्र ही है, और समता-स्थिति
मूल-गणना को बहुलताओं के साथ और तीक्ष्ण कर देती है। अंतिम शब्द प्रतिस्थापन $x
= \cos\theta$ का है: वह बहुपदों की दृढ़, विविक्त दुनिया को त्रिकोणमिति की
आवर्ती दुनिया में ले जाता है, जहाँ $T_n$ के मूल और चरम बिंदु केवल $\cos
n\theta$ की नियमित जाली हैं। विश्लेषण से उधार लिए दोनों तथ्य --- मध्यवर्ती
मान गुणधर्म (प्रश्न 14; \cref{ch:b1:continuity} में सिद्ध) और भीतरी चरम पर
अवकलज का लुप्त होना (प्रश्न 15; \cref{ch:b1:derivative} में सिद्ध) --- ठीक
वही औज़ार हैं जो आगे के अध्याय लौटाकर देंगे, और इस प्रकार चक्र पूरा हो जाता
है।
\end{solution}
